% =====================================================================
% tutuong.tex -- KHÔNG compile, KHÔNG \input từ main.tex.
% File ghi chú tư tưởng/quy tắc viết cho latex/paper/main.tex, để agent
% (Claude hoặc người) làm việc ở phiên sau đọc lại, tránh lặp lỗi đã sửa.
% Mỗi mục: [luật] + [lý do] + [cách phát hiện lại nếu tái phạm].
% =====================================================================

% ---------------------------------------------------------------------
% LUẬT 1 -- Không diễn giải lại cùng 1 luận điểm so sánh kỹ thuật
%           ở nhiều section khác nhau.
% ---------------------------------------------------------------------
% Lỗi đã xảy ra: 3 bài ghose2007validation, koliadis2006goalbpm,
% caballero2026alig từng được diễn giải chi tiết (cơ chế + điểm yếu)
% CẢ BA LẦN, gần như nguyên văn, ở Introduction (đoạn "State of the
% art"), Evaluation, và Related Work. Cụ thể câu về ghose2007validation
% ("reasoning is evaluated at a single point in time and cannot
% represent execution order") lặp lại gần như y hệt ở cả 3 nơi.
%
% Quy tắc áp dụng từ giờ: mỗi kỹ thuật đối sánh chỉ "sống" đầy đủ ở
% ĐÚNG MỘT section, những chỗ khác chỉ nêu tên + \cite + 1 mệnh đề
% ngắn, kèm \ref{} trỏ sang section giữ bản đầy đủ. Phân công cố định:
%   - Introduction: chỉ liệt kê hướng + cite, KHÔNG mô tả cơ chế từng
%     bài, KHÔNG lặp lại điểm yếu -- chỉ giữ đúng 1 câu "root
%     limitation" chung cho cả nhóm (đây là câu mở đầu, không lặp ở
%     đâu khác).
%   - Evaluation (Section~\ref{sect:evaluative}): giữ bản so sánh áp
%     dụng CỤ THỂ vào case study (Situation 1/2), vì đây là chỗ duy
%     nhất có ngữ cảnh cụ thể để so sánh có sức nặng.
%   - Related Work (Section~\ref{sect:relatedWork}): giữ bản ĐẦY ĐỦ
%     NHẤT, bao quát toàn bộ nhóm kỹ thuật (không chỉ 3 bài đại diện
%     nêu ở Introduction/Evaluation).
%
% Cách phát hiện lại: trước khi thêm câu văn mới nhắc tới 1 \cite{key}
% đã dùng ở nơi khác, chạy:
%   grep -n '\\cite{.*key.*}\|keyword-kỹ-thuật' main.tex
% để xem \cite{key} đã xuất hiện ở section nào, đọc lại câu cũ trước
% khi viết câu mới -- nếu ý gần giống, RÚT GỌN + trỏ \ref{} thay vì
% viết lại.
%
% NGOẠI LỆ: kỹ thuật "trỏ \ref{} sang section giữ bản đầy đủ" ở trên
% KHÔNG được áp dụng cho Introduction -- xem Luật 4, Introduction
% không được dùng \ref{} kiểu này dù để tránh trùng lặp.

% ---------------------------------------------------------------------
% LUẬT 2 -- Không dùng citation còn "TODO verify" trong references.bib
%           để khẳng định chắc chắn trong văn bản chính.
% ---------------------------------------------------------------------
% Lỗi đã phát hiện (CHƯA sửa, còn treo): groner2014vadl được dùng
% trong Related Work với câu khẳng định chắc "Gröner et al. extend the
% same technique to choreographies", nhưng entry trong references.bib
% còn nguyên dòng:
%   % TODO verify -- xác nhận đây có phải cùng bài với "paper_13"
%   trong SURVEY nội bộ hay không
% tức là nghi vấn "groner2014vadl có phải trùng với paper_13 (đoán
% nhầm tác giả) hay không" CHƯA được xác nhận qua DBLP, nhưng đã đưa
% thẳng vào bản thảo với tên tác giả cụ thể.
%
% Quy tắc áp dụng từ giờ: trước khi thêm/giữ 1 câu trích dẫn khẳng
% định chắc (tên tác giả, cơ chế cụ thể) cho 1 \cite{key} có dòng
% "% TODO verify" trong references.bib, PHẢI tự tra DBLP/venue gốc xác
% nhận trước, hoặc nếu chưa tra được thì viết dè dặt hơn (không nêu
% tên tác giả cụ thể) thay vì khẳng định chắc.
%
% Cách phát hiện lại: chạy
%   grep -B1 '^@' references.bib | grep -A1 'TODO verify'
% để liệt kê toàn bộ entry còn treo, đối chiếu với danh sách \cite{}
% thực sự dùng trong main.tex trước khi nộp bản thảo chính thức.
% n3_2019_reqdepgraph và n9_2025_goalsynthesis đã CHỦ ĐỘNG bị loại
% khỏi văn bản chính (chỉ còn trong .bib làm placeholder) vì tác giả
% chưa xác nhận đầy đủ -- đây là cách xử lý ĐÚNG, tiếp tục áp dụng
% cách này cho mọi entry còn TODO khi chưa xác nhận được.

% ---------------------------------------------------------------------
% LUẬT 3 -- Mỗi nội dung chỉ được viết vào ĐÚNG 1 trong 8 section,
%           theo đúng vai trò của section đó. Không ghi nhầm section.
% ---------------------------------------------------------------------
% main.tex hiện có đúng 8 section (theo thứ tự): Introduction,
% Background and Motivating Example, Overview of Our Approach, Tool
% Support, Evaluation, Discussion (kèm subsection Limitations),
% Related Work, Conclusion. Luật 1 đã xử lý 1 dạng lỗi cụ thể (lặp nội
% dung so sánh kỹ thuật); Luật 3 tổng quát hoá thành nguyên tắc: MỌI
% loại nội dung, không chỉ so sánh kỹ thuật, đều có đúng 1 "nhà" --
% viết vào đâu ngoài đó là ghi nhầm chỗ, dù không trùng lặp chữ nào.
%
% Vai trò từng section (được phép chứa gì -- KHÔNG được chứa gì):
%
%   1. Introduction (\label{sect:intro})
%      ĐƯỢC: bối cảnh bài toán, khoảng trống, RQ1/RQ2, tóm tắt 3 đóng
%      góp (1 câu/đóng góp), roadmap cuối section.
%      KHÔNG: mô tả cơ chế kỹ thuật của bài liên quan (chỉ nêu tên +
%      cite + \ref sang Evaluation/Related Work -- xem Luật 1); KHÔNG
%      chi tiết thuật toán/kiến trúc (đó là Overview); KHÔNG kết quả
%      thực nghiệm (đó là Evaluation/Tool Support).
%
%   2. Background and Motivating Example (\label{sect:motivation})
%      ĐƯỢC: case study MeetingScheduler, Situation 1/2/3, kiến thức
%      nền i*/BPMN/OCL-USE cần để đọc phần sau.
%      KHÔNG: so sánh với công trình liên quan; KHÔNG mô tả thuật toán
%      kiểm chứng (đó là Overview, chỉ được nhắc tên + \ref).
%
%   3. Overview of Our Approach (\label{sect:overview})
%      ĐƯỢC: ngôn ngữ ACL, luật dịch ACL sang USE/OCL, thuật toán
%      checkpoint-trace -- đây là phần đóng góp kỹ thuật cốt lõi.
%      KHÔNG: giao diện/plugin/tool (đó là Tool Support); KHÔNG kết
%      quả case study cụ thể ngoài ví dụ minh hoạ đã có sẵn ở
%      Background; KHÔNG so sánh với bài khác (đó là Evaluation).
%
%   4. Tool Support (\label{sect:toolSupport})
%      ĐƯỢC: kiến trúc plugin USE, 2 chế độ (batch/step-by-step), liệt
%      kê case study nào đã áp dụng (chỉ liệt kê, không phân tích kết
%      quả sâu).
%      KHÔNG: chi tiết thuật toán (đó là Overview); KHÔNG phân tích
%      trade-off/hạn chế (đó là Discussion).
%
%   5. Evaluation (\label{sect:evaluative})
%      ĐƯỢC: so sánh ÁP DỤNG CỤ THỂ với 3 kỹ thuật đại diện đã nêu tên
%      ở Introduction, gắn trực tiếp vào Situation 1/2.
%      KHÔNG: khảo sát đầy đủ 20 bài (đó là Related Work); KHÔNG bàn
%      trade-off thiết kế nội tại của phương pháp mình (đó là
%      Discussion, ví dụ "concrete trace vs exhaustive coverage" phải
%      nằm ở Discussion, không lặp lại ở đây).
%
%   6. Discussion (\label{sect:discussion}, có subsection Limitations)
%      ĐƯỢC: trade-off thiết kế (vd tách 3 loại failure, vết cụ thể vs
%      bao phủ toàn diện), hạn chế cụ thể của CHÍNH phương pháp này
%      (coverage, scenario quality, fail-fast masking, pairwise-only,
%      evaluation scope), câu hỏi mở.
%      KHÔNG: so sánh với bài khác (đó là Evaluation/Related Work);
%      KHÔNG mô tả case study (đó là Background/Tool Support).
%
%   7. Related Work (\label{sect:relatedWork})
%      ĐƯỢC: khảo sát ĐẦY ĐỦ NHẤT toàn bộ 5 nhóm kỹ thuật (khoảng 18
%      bài, xem latex/content/02-cong-trinh-lien-quan.md), mỗi nhóm:
%      cơ chế chung -> đại diện + cite -> điểm mạnh/yếu -> khác biệt
%      với phương pháp đề xuất. Đoạn tổng hợp cuối chốt lại RQ1/RQ2.
%      KHÔNG: nhắc lại nguyên văn câu đã dùng ở Introduction/Evaluation
%      cho 3 bài đại diện (xem Luật 1) -- vẫn được PHÉP nhắc lại tên 3
%      bài đó vì chúng cũng thuộc 1 trong 5 nhóm, nhưng câu văn phải
%      khác, không copy.
%
%   8. Conclusion (\label{sect:conclusion})
%      ĐƯỢC: tóm tắt 3 đóng góp, tóm tắt kết quả case study, tóm tắt
%      hạn chế đã nêu ở Discussion (1 câu, không giải thích lại chi
%      tiết), future work.
%      KHÔNG: lập luận/so sánh mới chưa xuất hiện ở section nào trước
%      đó -- Conclusion không phải chỗ để giới thiệu ý mới.
%
% Cách phát hiện lại: khi sắp thêm 1 câu/đoạn mới, tự hỏi "câu này
% thuộc về vai trò section nào trong 8 vai trò trên", rồi kiểm tra
% đang đứng trong \section nào bằng cách đọc ngược lên \label{sect:...}
% gần nhất phía trên vị trí sắp sửa/thêm. Nếu vai trò không khớp
% section hiện tại, chuyển đoạn đó sang đúng section, không viết tại
% chỗ đang đứng cho tiện.

% ---------------------------------------------------------------------
% LUẬT 4 -- Introduction phải đọc độc lập với toàn bài báo: không được
%           dùng \ref{} trỏ sang section khác, TRỪ đúng 1 đoạn cuối
%           cùng "The remainder of this paper is organized as
%           follows..." (đoạn bố cục paper).
% ---------------------------------------------------------------------
% Lý do (nguyên văn ý người dùng): Introduction phải đọc độc lập với
% toàn bài báo -- một người chỉ đọc Introduction (không đọc tiếp) vẫn
% phải hiểu trọn vẹn nội dung, không bị buộc phải lật sang section
% khác giữa chừng câu văn. Đoạn bố cục cuối Introduction là ngoại lệ
% DUY NHẤT, vì bản chất của nó chính là liệt kê section nào chứa gì.
%
% Lỗi đã xảy ra và đã sửa: khi áp dụng Luật 1 (tránh trùng lặp) cho
% đoạn "State of the art", đã lỡ thêm
%   "Sections~\ref{sect:evaluative} and~\ref{sect:relatedWork}
%   compare these directions in detail"
% vào giữa Introduction -- đúng kiểu \ref{} bị Luật 4 cấm. Đồng thời
% đoạn "Contributions" (4 câu "First/Second/Third/Fourth") cũng có sẵn
% 4 chỗ "(Section~\ref{...})" sau mỗi câu -- cũng vi phạm, dù không
% liên quan gì tới Luật 1/trùng lặp. Cả 2 chỗ đã được gỡ \ref{}, giữ
% nguyên câu chữ còn lại.
%
% Hệ quả với Luật 1: cách "trỏ \ref{} sang section giữ bản đầy đủ" để
% tránh trùng lặp KHÔNG dùng được cho Introduction. Cách đúng cho
% Introduction là giữ câu văn ở mức trừu tượng, đủ để tự đứng một
% mình, và để đoạn bố cục cuối cùng làm nhiệm vụ chỉ đường -- không
% chèn \ref{} rải rác vào các đoạn thân bài.
%
% Cách phát hiện lại: chạy
%   grep -n '\\ref{' main.tex | awk -F: '$1 >= <dòng-bắt-đầu-Introduction> && $1 <= <dòng-kết-thúc-Introduction>'
% hoặc đơn giản hơn: mở đúng vùng \section{Introduction}...\section{Background...},
% grep '\\ref{' trong vùng đó -- mọi kết quả ngoài đoạn "The remainder
% of this paper is organized as follows" đều là vi phạm cần gỡ.
% Lưu ý: \cite{} (trích dẫn tài liệu) KHÔNG bị cấm, luật này chỉ cấm
% \ref{sect:...}/\ref{sec:...} (tham chiếu chéo nội bộ bài báo).

% ---------------------------------------------------------------------
% LUẬT 5 -- \caption của hình phải ngắn gọn, không lan man.
% ---------------------------------------------------------------------
% Chỉ nêu tên/ý chính của hình bằng cụm ngắn, không viết lại thành câu
% mô tả hết mọi chi tiết đã thấy được trong hình.
% Sai:  "The ACL model of MeetingScheduler: roles Initiator/Organizer/
%        Secretary/Participant specializing MeetingParty, owned by
%        group MeetingUnit, with the Initiator--Participant and
%        Organizer--Secretary compatibility links and the Meeting/
%        PhoneContact entities."
% Đúng: "The ACL model of MeetingScheduler."
% CHƯA áp dụng cho main.tex -- hiện có khoảng 8 caption hình đang dài
% kiểu trên (chạy grep '\\caption{' main.tex để liệt kê).

% ---------------------------------------------------------------------
% LUẬT 6 -- Hạn chế tối đa dấu chấm phẩy (;) trong văn xuôi main.tex.
% ---------------------------------------------------------------------
% Lý do (nguyên văn ý người dùng): câu nối nhiều mệnh đề bằng ";" rất
% khó đọc với người, dù ngữ pháp đúng. Hiện main.tex đang có 54 dấu ";"
% trên 39 dòng (grep ';' main.tex).
%
% Quy tắc: ưu tiên tách thành 2 câu riêng (dấu chấm), hoặc dùng dấu hai
% chấm/dấu phẩy nếu mệnh đề sau chỉ bổ nghĩa/liệt kê ngắn. Chỉ giữ ";"
% khi nối 2 mệnh đề độc lập RẤT ngắn và có liên hệ nhân-quả/tương phản
% tức thời, không dùng để nhồi 3+ mệnh đề vào 1 câu.
% Không áp dụng cho itemize (mỗi \item kết bằng ";" theo latex-rules là
% quy ước khác, không liên quan) và không áp dụng cho ghi chú trong
% chính tutuong.tex.
%
% Cách phát hiện lại: chạy `grep -c ';' main.tex`, xem lại từng dòng có
% ";" -- câu nào có 2+ dấu ";" hoặc câu dài hơn ~30 từ thì gần như chắc
% cần tách.

% ---------------------------------------------------------------------
% Trạng thái tại thời điểm ghi chú này (để agent sau biết đã sửa gì)
% ---------------------------------------------------------------------
% - Đã sửa: cắt trùng lặp Introduction/Evaluation/Related Work theo
%   Luật 1 (xem lịch sử hội thoại / git diff main.tex quanh các dòng
%   "State of the art", Evaluation đoạn đầu, Related Work nhóm B).
% - Đã chuyển nội dung "business-irrational verdicts" (điểm yếu riêng
%   của caballero2026alig) từ Introduction (bị cắt) sang đúng chỗ ở
%   Related Work nhóm "Operational semantics along an execution
%   trace", để không mất nội dung khi rút gọn Introduction.
% - CHƯA sửa: rủi ro groner2014vadl ở Luật 2 -- vẫn cần tự tra DBLP
%   trước khi nộp bản chính thức.
% - Đã sửa: gỡ 2 chỗ \ref{} vi phạm Luật 4 trong Introduction (đoạn
%   "State of the art" và đoạn "Contributions").

% =====================================================================
% MỤC APPROACH -- Cách viết 1 subsection giới thiệu ngôn ngữ/metamodel
% mới (metamodel -> concrete syntax -> semantics), rút ra từ session
% viết lại sec:acl-language ("The ACL Language"). Dùng mục này để
% hướng dẫn agent khác viết các subsection tương tự (i*, BPMN...),
% không chỉ áp dụng riêng cho ACL.
% =====================================================================

% ---------------------------------------------------------------------
% A1 -- Thứ tự trình bày cố định: Metamodel -> Concrete syntax ->
%       Semantics. Không xen phần "dịch/convert sang công cụ khác"
%       (VD ACL-to-USE) vào giữa, trừ khi bài thực sự cần trình bày
%       bước dịch đó như một đóng góp riêng.
% ---------------------------------------------------------------------

% ---------------------------------------------------------------------
% A2 -- Phân biệt rõ 2 tầng: "schema" (mô hình tuân theo metamodel, vd
%       ACL) và "marking/instance" (1 configuration cụ thể lúc chạy,
%       vd AOL / System marking). Mỗi tầng có ĐÚNG 1 định nghĩa hình
%       thức (tuple); không viết công thức tuple cho tầng nào 2 lần.
%       Nếu tầng schema đã có hình UML metamodel làm nhiệm vụ đó rồi,
%       CÓ THỂ không cần công thức hoá tầng schema riêng -- gộp thẳng
%       vào định nghĩa tầng marking (vd: "An ACL model is a tuple
%       A=(...), instantiating Figure X: ... A system marking on A is
%       a tuple Sigma=(...): ...", một Definition duy nhất).
%
% Lý do: hình UML + công thức tuple mô tả cùng 1 thứ (schema) ở 2 dạng
% ký pháp khác nhau -- viết cả 2 đầy đủ là trùng lặp không cần thiết.

% ---------------------------------------------------------------------
% A3 -- Công thức tuple chỉ viết ĐẦY ĐỦ 1 LẦN trong toàn bài. Định
%       nghĩa khác cần dùng lại (vd BPMN model definition dùng lại ACL
%       model A) chỉ trỏ "$\mathcal{A}$ (Definition~\ref{def:...})",
%       KHÔNG liệt kê lại các thành phần "(E,R,G,Assoc,Compat)".
%
% Cách phát hiện lại: grep tên ký hiệu tuple (vd 'mathit{Assoc}') --
% nếu xuất hiện y hệt ở 2+ chỗ trong main.tex, chỉ giữ 1 chỗ định
% nghĩa đầy đủ, chỗ còn lại rút gọn thành tham chiếu.

% ---------------------------------------------------------------------
% A4 -- Không hình thức hoá (formalize) thành phần không được dùng lại
%       ở đâu khác trong bài. Trước khi thêm 1 ký hiệu/hàm vào công
%       thức (vd hàm cardinality $\mathit{card}$, thuộc tính
%       scope/direction của 1 relation), grep xem ký hiệu đó có xuất
%       hiện lại ở chỗ nào khác không -- nếu không, bỏ khỏi công thức,
%       chỉ nhắc bằng lời (nếu cần) ở phần mô tả metamodel.

% ---------------------------------------------------------------------
% A5 -- Không câu "rào" (hedging/filler). Cắt:
%       - câu roadmap kiểu "This section presents X, then Y, before Z";
%       - câu tự biện minh kiểu "the diagrammatic one suffices here,
%         since nothing in X needs...";
%       - câu diễn giải lại 1 ý đã nói bằng cách khác (kiểu paraphrase
%         không thêm thông tin mới).
% Chỉ thêm mệnh đề giải thích/rào đón khi người đọc THỰC SỰ cần chi
% tiết đó để hiểu đúng, không thêm cho "chắc ý".

% ---------------------------------------------------------------------
% A6 -- Hạn chế \texttt{} cho các từ lấy trực tiếp từ ví dụ minh hoạ
%       (tên role/entity cụ thể như Initiator, Meeting, PhoneContact...)
%       -- viết bằng văn xuôi thường ("an initiator, an organizer, an
%       optional secretary, and several participants"), không code-font.
%       CHỈ giữ \texttt{} cho:
%       (i) đuôi file/tên ngôn ngữ (.acl, .aol);
%       (ii) định danh literal lấy nguyên văn từ 1 Listing cụ thể
%            (org1, sec1, a2 -- vì đó là token thật trong code, không
%            phải tên khái niệm);
%       (iii) tên lớp/khái niệm metamodel còn được DÙNG LẠI về sau
%            trong bài bằng công thức hình thức (vd Owner, Compatibility
%            nếu công thức closed-world còn nhắc tới).

% ---------------------------------------------------------------------
% A7 -- Không lặp cùng 1 nội dung dưới 2 dạng thể hiện (concrete syntax
%       dạng chữ/Listing + dạng hình/Figure) nếu 2 cái show cùng 1 thứ.
%       Giữ 1, ưu tiên HÌNH nếu đã có hình tương ứng; chỉ dùng Listing
%       khi KHÔNG có hình minh hoạ tương ứng cho đúng artifact đó.

% ---------------------------------------------------------------------
% A8 -- Khi cắt 1 khái niệm/ví dụ khỏi bản ngắn (conference version)
%       nhưng có thể cần lại cho bản dài (journal/long version), KHÔNG
%       xoá hẳn -- chuyển sang file appendix-<ten>.tex riêng, theo đúng
%       convention đã có (xem appendix-algorithms.tex,
%       appendix-aol-example.tex): comment đầu file nêu rõ (a) lý do
%       cắt, (b) cách dùng lại (\input hoặc copy đoạn nào vào đâu),
%       (c) những \ref{}/\label{} ở main.tex cần khôi phục nếu ghép
%       lại. Đồng thời phải tự tìm và sửa MỌI \ref{} trong main.tex
%       từng trỏ tới nội dung vừa cắt (kể cả ở section khác), không để
%       lại tham chiếu treo (hiện "??" khi compile).

% ---------------------------------------------------------------------
% A9 -- Bảng (table) chỉ giữ lại nếu nó nêu được điều câu văn ngay
%       trước KHÔNG thể nêu gọn bằng lời (vd dữ liệu nhiều chiều, số
%       liệu). Nếu bảng chỉ lặp lại y hệt điều câu văn đã nói (vd bảng
%       3 dòng "role pair / declared? / outcome" trong khi câu văn đã
%       nêu đủ cả 3 dòng đó rồi), CẮT bảng.

% ---------------------------------------------------------------------
% A10 -- Mở đầu đoạn giới thiệu 1 khái niệm/metamodel bằng LÝ DO/vấn đề
%        cần giải quyết trước, rồi mới tới cấu trúc. Không mở thẳng
%        bằng "Figure X: class A specializes class B, class C có thuộc
%        tính..." mà không có câu ngữ cảnh động lực đứng trước.
% Sai:  "Figure~\ref{fig:x} gives Y's abstract syntax: A specializes B..."
% Đúng: "An organization needs more than isolated roles: a unit can
%        itself contain sub-units... Figure~\ref{fig:x} gives Y's
%        abstract syntax for exactly this: ..."

% ---------------------------------------------------------------------
% A11 -- Khi metamodel mới mở rộng từ 1 metamodel/ngôn ngữ có sẵn (vd
%        UML class diagram đã có sẵn trong USE), nói rõ CÁI GÌ KẾ THỪA
%        SẴN CÓ và CÁI GÌ LÀ THÊM MỚI riêng của ngôn ngữ đang giới
%        thiệu (kiểu "extend X with Y to support Z" -- xem cách bài
%        UDML viết AS của họ). Không liệt kê phẳng mọi lớp như thể tất
%        cả đều ngang hàng mới, người đọc không phân biệt được đâu là
%        máy sẵn có, đâu là đóng góp riêng.

% ---------------------------------------------------------------------
% A12 -- Khi mô tả 1 class diagram bằng lời, đi theo lối "dẫn từ gốc
%        xuống lá" (root-to-leaf), như đang chỉ tay lướt từng vùng của
%        hình theo thứ tự: container/gốc trước -> lớp trừu tượng dùng
%        chung -> các lớp con -> quan hệ/lớp ngoại vi -> phần mở rộng.
%        KHÔNG liệt kê theo kiểu "A specializes B; C specializes D;
%        E is..." phẳng theo alphabet hay theo nhóm loại (dễ đọc như
%        bản kê khai, không như một lời giải thích).

% ---------------------------------------------------------------------
% A13 -- Mỗi đoạn/khối (Metamodel, Concrete syntax, Semantics...) phải
%        MỞ ĐẦU bằng 1 câu nối tường minh với khối ngay trước (vd
%        "X instantiates this metamodel directly: ...", "What such a
%        schema actually admits at runtime... is not fixed by the
%        schema itself; it takes a separate notion... to pin down
%        one concrete configuration."). KHÔNG chỉ đặt các khối cạnh
%        nhau dưới nhãn in nghiêng rời rạc (`\emph{Metamodel.}`,
%        `\emph{Concrete syntax.}`, `\emph{Semantics.}`...) -- nhãn
%        kiểu đó dễ gây cảm giác "liệt kê từng mục" thay vì 1 mạch văn
%        liên tục có lập luận. Có thể bỏ hẳn các nhãn in nghiêng này
%        nếu câu mở đầu mỗi đoạn đã đủ rõ đang chuyển sang phần nào.

% ---------------------------------------------------------------------
% A14 -- Definition hình thức có nhiều thành phần kiểu tuple nên trình
%        bày bằng itemize (mỗi thành phần 1 dòng), không nhồi hết vào
%        1 câu văn nối bằng dấu chấm phẩy (dễ đọc "cụt", khó dò từng
%        thành phần). Xem cách bài UDML viết Definition cho AS/CS.

% ---------------------------------------------------------------------
% A15 -- `\begin{definition}[...]` đã tự đóng vai trò báo hiệu "đây là
%        phần ngữ nghĩa hình thức" -- KHÔNG cần thêm nhãn
%        `\emph{Semantics.}` ngay trước nó (2 tín hiệu cho cùng 1 việc
%        là thừa). Nếu cần dẫn nhập trước Definition, dẫn bằng câu nêu
%        CÂU HỎI mà Definition trả lời (vd "it takes a separate
%        notion... to pin down one concrete configuration"), không
%        phải câu thông báo tên phần.

% ---------------------------------------------------------------------
% A16 -- Khi giải thích cơ chế kiểm tra/validity có nhiều trường hợp,
%        nêu rõ "có mấy cách" trước rồi mới liệt kê từng cách (vd
%        "Validity is enforced two ways: cardinality and typing follow
%        directly from A's structure, while compatibility is evaluated
%        closed-world..."), thay vì để người đọc tự suy ra đang có bao
%        nhiêu cơ chế trong đoạn văn.

% ---------------------------------------------------------------------
% A17 -- Trước khi chốt 1 cụm đoạn (metamodel + definition + đoạn ngay
%        sau definition), rà lại: mỗi Ý chỉ được xuất hiện ĐÚNG 1 LẦN
%        trong CẢ CỤM, không chỉ trong 1 câu/1 đoạn riêng lẻ. 1 ý bị
%        lặp giữa đoạn Metamodel + Definition + đoạn giải thích ngay
%        sau Definition vẫn tính là lặp dù nằm ở 3 vị trí khác nhau
%        (vd "compatibility chỉ nối 2 role" từng bị nói tới 3 lần liền
%        nhau ở 3 đoạn khác nhau -- chỉ giữ đúng 1 chỗ, ưu tiên chỗ
%        cần ký hiệu đó nhất để dùng ngay sau).
%
% Cách phát hiện lại: đọc liền 3-4 đoạn liên tiếp thành 1 khối, tự hỏi
% "câu này có đang nói lại đúng ý câu nào ở đoạn trước không" -- nếu
% có và không thêm thông tin mới, cắt.

% =====================================================================
% LUẬT 6 -- Formalize attribute values bằng link tới object có sẵn
%           (primitive value), KHÔNG cần thêm hàm val riêng.
% =====================================================================
% Quyết định (Definition [System marking], sec:acl-language, def:system-
% marking): thay vì $\Sigma = (Agent, O, type, val, link)$ với val gán
% giá trị thuộc tính riêng, ta coi mỗi primitive type (Boolean,
% Integer, String...) là 1 "entity có sẵn" nằm trong $E$ (do class-
% diagram/metamodel cấp sẵn -- xem Definition [ACL model], def:acl-
% model -- người dùng ACL không cần tự khai báo Boolean), và giá trị
% của nó (true, false, mỗi số nguyên, mỗi string...) là 1 object có
% sẵn trong $O$ với $type(\texttt{true}) = type(\texttt{false}) =
% \mathit{Boolean}$. Gán thuộc tính (vd m1.detailsDecided = false) khi
% đó chỉ là 1 link bình thường $(m1, \texttt{false}) \in link$, không
% cần hàm val riêng.
%
% Kết quả: $\Sigma$ rút gọn còn 4 phần $(Agent, O, type, link)$ thay vì
% 5, và $type: O \to E\cup R\cup G$ vẫn chỉ có đúng 3 loại đích (không
% cần mở rộng thêm 1 miền DataType/primitive-value riêng).
%
% Lý do giữ cách này:
% (i) đơn giản hơn -- val không được dùng lại ở công thức nào khác
%     trong toàn bài (đã grep 'mathit{val}' kiểm tra trước khi bỏ), nên
%     bỏ nó không mất thông tin ở chỗ khác;
% (ii) link đã có sẵn kiểu $(Agent\cup O)\times O$ nên không cần mở
%      rộng type signature gì thêm để chứa việc gán thuộc tính;
% (iii) nhất quán với cách metamodel (Definition [ACL model]) đã đặt
%      DataType/PrimitiveType làm anh em với Class dưới Classifier --
%      coi Boolean là 1 Entity có sẵn tránh phải thêm 1 tầng khái niệm
%      mới chỉ để xử lý riêng giá trị nguyên thủy.
%
% Điều cần nhớ nếu sau này viết lại công thức này: ĐỪNG thêm lại 1 hàm
% val riêng trừ khi thực sự có 1 công thức khác trong bài cần thao tác
% trực tiếp trên "giá trị" (khác bản chất với "object") -- nếu gặp
% trường hợp đó, cân nhắc kỹ trước khi revert lại cách cũ, vì bản thân
% cách val riêng KHÔNG sai, chỉ là thừa so với những gì bài hiện dùng.
%
% Cách phát hiện lại: grep 'mathit{val}' trong main.tex -- nếu thấy nó
% xuất hiện trở lại trong Definition [System marking], đối chiếu lại
% note này trước khi quyết định giữ hay bỏ.
